\suggestions
 
A partir de los resultados obtenidos y las limitaciones identificadas durante el desarrollo, se
formulan las siguientes recomendaciones para trabajos futuros:
 
\begin{enumerate}
 
  \item \textbf{Incorporar soporte para subida paralela de fragmentos.} La implementación actual
  transfiere los fragmentos de forma secuencial para simplificar la gestión del estado y la
  detección de errores por fragmento. Una evolución natural del módulo consistiría en habilitar la
  transferencia concurrente de múltiples fragmentos mediante un pool de promesas controlado,
  reduciendo el tiempo total de subida en conexiones de alta capacidad sin comprometer la
  consistencia del estado persistido. Esta mejora requeriría ajustar la lógica de reconciliación
  del progreso tanto en el cliente como en el endpoint \texttt{complete} del backend.
 
  \item \textbf{Sustituir el mecanismo de \textit{polling} por notificaciones en tiempo real
  mediante WebSockets.} El módulo implementado consulta periódicamente el endpoint
  \texttt{/verificar\_cola/} cada tres segundos para conocer el estado del procesamiento, lo que
  introduce latencia artificial en la actualización de la interfaz y genera tráfico HTTP innecesario
  durante períodos de espera prolongados. Se recomienda reemplazar este mecanismo por un canal
  WebSocket ---integrable en el ecosistema Django mediante Django Channels--- que permita al
  \textit{worker} de Celery notificar al cliente de forma inmediata al completar el procesamiento.
  Esta mejora eliminaría el retardo perceptible entre la finalización real de la tarea y su
  reflejo en la interfaz, y reduciría la carga de solicitudes sobre el servidor en escenarios con
  múltiples usuarios procesando contenido de forma simultánea.
 
\end{enumerate}